\section{Introduction}
\label{introduction}

\subsection{Background and Motivation}
\label{sec:intro-background}

Natural-language requirements are the main interface between stakeholder intent and software design.
However, requirements are often written without explicit architectural assumptions, component responsibilities, quality-attribute trade-offs, or design constraints.
This gap makes it difficult for intelligent software systems to move from requirement understanding to architecture-aware design decisions.
Recent requirements engineering research has also emphasized the central role of requirements in intelligent software development~\cite{robinsonRequirementsAreAll2025} and the persistent difficulty of writing precise yet stakeholder-readable requirement specifications~\cite{darifControlledNaturalLanguage2026}.

\subsection{Challenges of Requirement-to-Architecture Alignment}
\label{sec:intro-challenges}

This paper focuses on the requirement-to-architecture alignment problem.
The key challenges are: identifying architectural concerns from informal requirements, mapping requirements to components and responsibilities, aligning quality attributes with architecture constraints, and preserving traceability between requirement statements and design decisions.

\subsection{Multi-Agent Architecture Reasoning}
\label{sec:intro-multi-agent}

We investigate a lightweight multi-agent reasoning workflow for architecture-aware requirements engineering.
The workflow separates semantic requirement analysis, architecture reasoning, and review-oriented consistency checking, while keeping human feedback in the loop for ambiguous or high-impact mapping decisions.

\subsection{Contributions}
\label{sec:intro-contributions}

The paper is structured around the following research contributions:
\begin{itemize}[leftmargin=*]
  \item A clear formulation of requirement-to-architecture alignment for intelligent software systems.
  \item A knowledge-guided multi-agent workflow for extracting architectural concerns and constructing requirement-to-architecture mappings.
  \item A traceability-oriented mapping method covering component identification, responsibility allocation, and architecture constraint alignment.
  \item An empirical evaluation plan for assessing alignment accuracy, mapping completeness, traceability quality, and architecture consistency.
\end{itemize}
